\documentclass[10pt]{amsart}

\usepackage[margin=1in]{geometry}
\usepackage[T1]{fontenc}
\usepackage{lmodern}
\usepackage{amsmath,amssymb,amsthm,mathtools}
\usepackage{microtype}
\usepackage{url}
\usepackage{enumitem}
\usepackage{alltt}
\usepackage{booktabs}
\usepackage[dvipsnames]{xcolor}
\usepackage[colorlinks=true,linkcolor=MidnightBlue,citecolor=OliveGreen,urlcolor=MidnightBlue]{hyperref}
\usepackage[capitalize,nameinlink]{cleveref}

\setlist{leftmargin=*,itemsep=2pt,topsep=4pt}
\emergencystretch=2.5em

\theoremstyle{plain}
\newtheorem{theorem}{Theorem}[section]
\newtheorem{lemma}[theorem]{Lemma}
\newtheorem{proposition}[theorem]{Proposition}
\newtheorem{corollary}[theorem]{Corollary}
\newtheorem{importedtheorem}[theorem]{Imported Theorem}
\theoremstyle{definition}
\newtheorem{definition}[theorem]{Definition}
\newtheorem{example}[theorem]{Example}
\theoremstyle{remark}
\newtheorem{remark}[theorem]{Remark}
\crefname{importedtheorem}{Imported Theorem}{Imported Theorems}
\numberwithin{equation}{section}

\newcommand{\A}{\textup{[A]}}
\newcommand{\Oa}{\textup{[O]}}
\newcommand{\Z}{\mathbb{Z}}
\newcommand{\N}{\mathbb{N}}
\newcommand{\hF}{h_{\mathrm{form}}}
\newcommand{\val}{\mathrm{val}}
\newcommand{\rate}{2^{O(n/\sqrt{\log n})}}
% Lean names: typewriter, with line breaks allowed after '.' and '_'.
\DeclareUrlCommand\lean{\urlstyle{tt}\def\UrlBreaks{\do\.\do\_}\def\UrlBigBreaks{}}
% Repository names: typewriter, with line breaks allowed after '/' and '-'.
\DeclareUrlCommand\repo{\urlstyle{tt}\def\UrlBreaks{\do\/\do\-}\def\UrlBigBreaks{}}

\begin{document}

\title[Class numbers, square roots and non-residues in $2^{O(n/\sqrt{\log n})}$]{Class numbers, square roots and quadratic non-residues\\ in time $2^{O(n/\sqrt{\log n})}$, assuming the new $\#$SAT algorithms}

\author{Vaibhav Tulsyan}

\date{November 2026}

\subjclass[2020]{11Y16, 11Y40, 11A15, 68Q25, 03B35}
\keywords{Class number, binary quadratic forms, quadratic non-residue, square roots modulo $m$, divisor summatory function, Gauss circle problem, exact counting, Exponential Time Hypothesis, Lean~4, formal verification}

\begin{abstract}
Basic problems of computational number theory have proven deterministic bounds that are fixed powers of
the input. The class number $h(D)$ of an imaginary quadratic discriminant takes $|D|^{1/2+o(1)}$ steps
unconditionally and $|D|^{1/5+o(1)}$ under the generalised Riemann hypothesis (GRH); the least quadratic
non-residue modulo a prime $p$ takes $p^{1/(4\sqrt e)+o(1)}$ by Burgess's 1957 bound, polynomial only
under GRH; exact divisor sums $\sum_{m\le x}\tau(m)$ take $x^{1/3+o(1)}$. However, none of these bounds is
of the form $X^{o(1)}$, even conditionally. We present deterministic algorithms in time
$2^{O(n/\sqrt{\log n})}$, where $n$ is the input bit length, for six such problems: the number of reduced
forms of discriminant $-N$, for fundamental discriminants $h(-N)$; the least square root of $a$
modulo any $m$; the least quadratic non-residue modulo a prime; and $\sum_{m\le x}\tau(m)$,
$\sum_{m\le x}\sigma(m)$ and the Gauss circle count. Our approach writes each answer as a count of solutions
of equations with one or two products and evaluates it with Anthropic's counting engine through one
network library: products by a two-dimensional Sch\"onhage--Strassen network, sums by signed carry
chains, inequalities by comparison chains, with $O(n)$ shared bits. We prove the seven resulting theorems
in Lean~4 (for example \texttt{hForm\_TM\_rate}; axioms \texttt{propext}, \texttt{Classical.choice},
\texttt{Quot.sound}): a finite-alphabet Mathlib Turing machine outputs the answer within
$\lfloor C\cdot2^{Cn/\sqrt{\log_2(n+2)}}\rfloor$ steps on every input. Lean proves the count of Gauss's
reduced forms; its equality with the ideal class number of $\mathbb{Q}(\sqrt{-N})$ is classical and not
formalised. The constants are astronomical, the machines use exponential space, and every result rests
on Anthropic's Lean development of its claimed refutation of the Exponential Time Hypothesis, which is
under review.
\end{abstract}

\maketitle

% =====================================================================
\section{Introduction}\label{sec:intro}
% =====================================================================

Computing the class number $h(D)$ of an imaginary quadratic discriminant $D=-N$ is one of the oldest
problems of computational number theory. Gauss's reduction theory gives the simplest algorithm: $h(D)$ is
the number of reduced forms $ax^2+bxy+cy^2$ of discriminant $D$, and every such form has $a\le\sqrt{N/3}$,
so listing them takes $N^{1/2+o(1)}$ steps. To our knowledge this is still the best proven
deterministic bound without hypotheses~\cite[\S5]{Len92}. Under GRH, Shanks's baby-step giant-step
method, analysed by Lenstra~\cite{Len82} and Schoof~\cite[\S3]{Sch82}, takes $N^{1/5+o(1)}$, and the
subexponential algorithm of Hafner and McCurley~\cite{HM89} is randomised and also rests on GRH~\cite[\S5]{Len92}.
Two neighbouring problems have stood still for as long. The least quadratic non-residue modulo a prime
$p$ is below $p^{1/(4\sqrt e)+\varepsilon}$ by Burgess~\cite{Bur57}, so trial search takes
$p^{0.1517}$ steps, and this exponent has not been improved since 1957; only under GRH is the
non-residue below $O(\log^2p)$ (Ankeny~\cite{Ank52}, Bach~\cite{Bac90}). Exact divisor sums
$\sum_{m\le x}\tau(m)$ take $x^{1/3+o(1)}$ steps~\cite{Sla12,TCH12}.

In October 2026 OpenAI~\cite{OpenAI26} and Anthropic~\cite{Anthropic26} (below \Oa{} and \A) each
claimed a deterministic algorithm that counts the satisfying assignments of a $k$-CNF in subexponential
time, refuting the Exponential Time Hypothesis; both claims come with Lean~4 developments and are under
review. Rudich and Rousseau~\cite{RR26} used the engines to factor in $N^{o(1)}$, Shunia~\cite{Shu26}
proved the rate $\rate$ on paper, and our factoring note~\cite{TulFac26} proved it in Lean on \A's engine.
However, the bounds above are all fixed powers of the input, even under GRH for $h(D)$. We are unaware of
any derivation of them from the new counting engines; our searches are described in \cref{sec:related}.

\subsection{Counting problems with one or two products}
We consider problems whose answer is, or is found from, the number of integer solutions of a fixed
system of polynomial equations and inequalities in $O(1)$ unknowns of $O(n)$ bits, in which each
equation has at most two products. The class number is such a count:
$h(-N)$ counts the triples $(a,b,c)$ with $b^2+N=4ac$ and $|b|\le a\le c$, up to a tie-break.

\begin{example}[Running example]\label{ex:23}
Take $N=23$. The reduced forms of discriminant $-23$ are $(a,b,c)=(1,1,6)$, $(2,1,3)$ and $(2,-1,3)$,
so $h(-23)=3$. Our machine asks two counts. Query~1 counts the triples of naturals with $b^2+23=4ac$ and
$b\le a\le c$; they are $(1,1,6)$ and $(2,1,3)$. Query~2 counts the triples with $b^2+23=4ac$ and
$0<b<a<c$; it is $(2,1,3)$ alone, which stands for the form $(2,-1,3)$. The answer is $2+1=3$.
\end{example}

\subsection{The naive encodings fail; the principle}
At first one might reduce each count to $\#$SAT and apply the engines' headline theorems. However,
Tseitin's transformation of a multiplication circuit has $\Omega(n\log n)$ clauses for every known
circuit, and \A's clause bound then gives $2^{O(n\sqrt{\log n})}$, worse than the classical bounds. Our
factoring note~\cite{TulFac26} showed how to avoid this for one equation $pq=N$: \A's network model
charges $2^{e\cdot U}$ for $U$ bits shared between nodes, with $e=O(1/\sqrt{\log n})$, and a
two-dimensional Sch\"onhage--Strassen transform multiplies with $O(n)$ shared bits. Writing a new
network for each problem would cost thousands of Lean lines each. Instead, we rely on a different
principle: \emph{a fixed system of equations with few products has a network assembled from three kinds of
block (products, carry chains and comparison chains), each cutting $O(n)$ bits, so one library proved once
serves every problem.} The per-problem work is a data value, the description, together with a
counting identity in pure Mathlib.

\subsection{What we prove}
In this paper we build that library, which we call ArithNet, write six descriptions, and prove in Lean
that one Turing machine per problem computes the answer. All six theorems have the same form.

\begin{theorem}[\cref{sec:statement}; machine-checked]\label{thm:main}
For each of the following functions $F$ there are a Turing machine $M$ in Mathlib's \lean{TM2} model,
with finite alphabets on all stacks, and a constant $C\in\N$, such that on every input word $w$ of
length $n$ (as specified) $M$ halts within $\lfloor C\cdot2^{Cn/\sqrt{\log_2(n+2)}}\rfloor$ steps with
output \lean{encodeNat}~$F$:
\begin{enumerate}[label=(\roman*)]
\item \lean{hForm_TM_rate}: $w=$ \lean{encodeNat}~$N$, $F=\hF(-N)$, the number of reduced forms of
discriminant $-N$ (\cref{def:hform});
\item \lean{hFormPrim_TM_rate}: the same machine, for $-N$ a fundamental discriminant, outputs the number of
primitive reduced forms, the form class number $h(-N)$;
\item \lean{sqrtMod_TM_rate}: $w=$ \lean{encodePair}~$m$~$a$, $F$ the least $y<m$ with $y^2\equiv a\pmod m$,
and $m$ if there is none;
\item \lean{lqnr_TM_rate}: $w=$ \lean{encodeNat}~$p$ with $p$ prime, $F$ the least quadratic non-residue
modulo $p$ ($p$ if there is none, i.e.\ $p=2$); on every other $p$ the machine halts within the
same bound with some $y\le p$;
\item \lean{divisorSum_TM_rate}, \lean{sigmaSum_TM_rate}, \lean{circle_TM_rate}: $w=$ \lean{encodeNat}~$x$,
$F=\sum_{m\le x}\tau(m)$, $\sum_{m\le x}\sigma(m)$, and $\#\{(u,v)\in\Z^2:u^2+v^2\le x\}$.
\end{enumerate}
\end{theorem}

\noindent Since $n=\log_2N+O(1)$, the bound in (i) is $N^{O(1/\sqrt{\log\log N})}=|D|^{o(1)}$, against
$|D|^{1/2+o(1)}$ unconditionally and $|D|^{1/5+o(1)}$ under GRH. In (iv) it is $p^{o(1)}$, against
$p^{0.152}$; in (v) it is $x^{o(1)}$, against $x^{1/3+o(1)}$. To our knowledge, no deterministic
algorithm with a bound of the form $X^{o(1)}$ was known for (i), (iii) or (v), even under GRH.

\subsection{Contributions}
Our contributions are:
\begin{itemize}
\item We introduce ArithNet, a library that turns a \emph{description} (variables, products of a variable
by a variable or a constant, linear equations with coefficient mass at most $32$, comparisons, bit
patterns) into a parsimonious network, and prove in Lean that its value is the number of solutions,
that its shared bits total $O(n)$ and that each of its $g$ nodes holds $O(n/g+n^{3/4})$ bits
(\cref{sec:arithnet}).
\item We present generic counted prefix-search and one/two-query machines that turn any polynomial-time
description reader into a Turing machine running in $\rate$ steps (\cref{sec:machines}).
\item We give the six descriptions and their counting identities, including the two-query split of
reduced forms and an exact formula for non-residues with a prefix from a count of squares
(\cref{sec:instances}).
\item We prove the six theorems of \cref{thm:main} in Lean, and state precisely what the class-number
theorem counts and what remains classical (\cref{sec:honesty,sec:statement,sec:formal}).
\end{itemize}
The ingredients (reduction theory, prefix search, the transform network) are classical or from our
factoring note; what we add is the library, its proofs, and the descriptions.

\subsection{Status of each claim}
We separate four kinds of claim. \emph{Machine-checked in Lean}: the seven theorems of \cref{thm:main},
with the axioms \lean{propext}, \lean{Classical.choice} and \lean{Quot.sound} only, and every lemma whose
Lean name we give. \emph{Proved on paper only}: the equality of the form class number with the ideal class
number of $\mathbb{Q}(\sqrt{-N})$ (classical; Dirichlet and Dedekind, see~\cite[Thm.~7.7]{Cox13}); the
decomposition of $\hF$ for non-fundamental discriminants in \cref{rem:nonfund}; the discussion of real
quadratic fields; and the comparisons with earlier bounds. \emph{Assumed}: nothing beyond the three axioms;
\A's results enter as Lean theorems of its development, which the kernel checks with ours.
\emph{Checked numerically}: the running examples and the decomposition of \cref{rem:nonfund} for $N<400$, by
a short exact computation that is not part of the proof; some of the small values, such as the split
$2+1$ of \cref{ex:23} and $\mathrm{lqnr}(7)=3$ of \cref{ex:qnr}, are also checked in Lean by \lean{decide}.

\subsection*{Outline}
\Cref{sec:engine} recalls \A's counting theorem. \Cref{sec:arithnet} builds ArithNet and
\cref{sec:machines} the machines. \Cref{sec:instances} gives the descriptions, and \cref{sec:honesty}
says what the class-number theorem counts. \Cref{sec:rate,sec:statement,sec:formal} cover the exponent,
the statements and the Lean development, and \cref{sec:related,sec:limitations} related work and
limitations.

% =====================================================================
\section{The engine counts accepted labellings}\label{sec:engine}
% =====================================================================

In this section we recall the one theorem of \A{} that we compose with, in the form used by our
factoring note~\cite[\S2]{TulFac26}. We use \A's Lean development at commit \texttt{16a5ce4}
(\texttt{anthropics/formal-math}, folder \texttt{eth}, Lean v4.33.1, Mathlib \texttt{0df444a}).

A \emph{network on $[g]$ with parameters $P,B,\tau$} has, for every node $a$, settings of at most $P$
bits and a weight $W_a$, and for every pair of nodes an interaction that factors through $c_{a,b}$
shared bits, with $\sum_bc_{a,b}\le P$; its value is $Z=\sum_z\prod_aW_a(z_a)\prod_{\{a,b\}}\Psi_{a,b}$ and
its total is $U=\sum_{a\ne b}c_{a,b}$. We use only \emph{label systems}: items (wires) placed on pairs of
nodes, each carrying a $W$-bit label seen by both nodes, with $W_a$ a $0/1$ table. The value is then the
number of labellings accepted by every table.

\begin{importedtheorem}[\A, Lemma~8(i) and Theorem~A(b); see~{\cite[Imported Thm.~2.2]{TulFac26}}]\label{imp:engine}
Let $K=K_\ell\ge2$ be a level of \A's schedule, $g=3K^8$, and $e_\ell\le(3x_0+6)/\sqrt{\log_2K}$ for an
absolute $x_0$. A network on $[g]$ with parameters $P,B,\tau$ and total $U$ can be evaluated within
$\bigl(K^g+2^{e_\ell K^2L}(K(L+B'))^{C_0K^2}2^{2gP/K}\tau\bigr)(g(P+B))^{C_1}$ operations, where
$L=\lceil(U+gP/K)/K^2\rceil$, $B'=\lceil2gP/K\rceil+g^2B$ and $C_0,C_1$ are absolute.
\end{importedtheorem}

\noindent In Lean this is \lean{ExplicitCost.groupingExplicit} together with the clique machine
of the counted recursion (\lean{CountedRoute.onTheCountedRecursion}); \A's host lemma \lean{CountKSat.exists_tmW} turns the
step-by-step machine into a Mathlib Turing machine with polynomial overhead. So a network family with
$U=O(n)$ and $gP=O(n)$ is evaluated in $\rate$ steps at the level $K^{40}\lesssim n$ of our factoring note.

% =====================================================================
\section{ArithNet: one network library for few-product systems}\label{sec:arithnet}
% =====================================================================

In this section we define descriptions and the count they denote, build the network of a description,
and state its value and parameter theorems. The definitions are pure Mathlib (module \lean{ArithDesc});
the network is in the bridge modules \lean{ArithParams} and \lean{ArithNet*}.

\subsection{Descriptions}
A description is plain data:
\begin{alltt}\footnotesize
inductive Operand | var (v : \(\mathbb{N}\)) | const (limbs : List \(\mathbb{N}\))
structure ProdD where u : \(\mathbb{N}\); w : Operand            -- the product x_u * w
inductive Src | prod (p : \(\mathbb{N}\)) | lin (v : \(\mathbb{N}\)) | const (limbs : List \(\mathbb{N}\))
structure Term where coeff : \(\mathbb{N}\); src : Src
structure ChainD where lhs : List Term; rhs : List Term
structure CmpD where v : \(\mathbb{N}\); rhs : Operand; rel : CmpRel   -- lt | le | eq | ge | gt
structure Desc where nV : \(\mathbb{N}\); pat : List (List (Option Bool));
  prods : List ProdD; chains : List ChainD; cmps : List CmpD
\end{alltt}
Constants are lists of $b$-bit limbs, with value $\sum_il_i2^{bi}$. A chain asserts
$\sum_{\mathrm{lhs}}\mathrm{coeff}\cdot\mathrm{value}=\sum_{\mathrm{rhs}}\mathrm{coeff}\cdot\mathrm{value}$,
where the value of a term is that of a product, a variable or a constant. A pattern fixes some bits of a
variable (\lean{none} is a free bit). \lean{Desc.Sat}~$b$~$D$~$x$ holds when every pattern matches, every
comparison holds and every chain balances, and the count of $D$ at size $n$ is
\[
  \mathrm{count}_n(D)=\#\bigl\{x\in[0,2^{n_B})^{\mathrm{nV}}:\ \texttt{Desc.Sat}\ b\ D\ x\bigr\},
\]
with $b$ the limb width and $n_B=\ell b>n$ the box width of~\cite[\S4]{TulFac26}. A description is well formed
(\lean{Desc.WFB}) when its indices are in range, its constants have at most $\ell$ limbs below $2^b$, its
patterns have length at most $n_B$, each side of each chain has coefficient mass at most $32$, and no
product occurs in two chains.

\subsection{Wires, vertices and checks}
Every wire carries a label of $W'=J+8$ bits, where $J=2^{\lfloor\nu/2\rfloor+7}$, $\nu=\lfloor\log_2(n+2)\rfloor$,
is the transform modulus exponent of~\cite{TulFac26}.
\begin{itemize}
\item \emph{Variables.} Each variable has $\ell$ limb wires, one copy per chain in which it occurs linearly,
and one copy per product side in which it is an operand.
\item \emph{Products.} Each product $x_u\cdot w$ gets its own two-dimensional Sch\"onhage--Strassen pipeline
over $\Z/(2^J+1)$: row and column transforms of both operands, pointwise products, inverse columns and
inverse rows, ending in the $2\ell$ coefficients of the product (the stages of~\cite[\S4]{TulFac26}).
A constant operand enters through limb wires whose values the cells fix.
\item \emph{Chains: signed carries.} Chain $c$ has a carry vertex at each position $k$, which reads the
coefficient digits $d^{\mathrm{L}}_k,d^{\mathrm{R}}_k$ of its two sides (sums of coefficient times product
coefficient, limb or constant limb) and the carry wires, and checks the single equation
\begin{equation}\label{eq:carry}
  R_{k+1}\cdot2^b+d^{\mathrm{R}}_k+\mathrm{OFF}=d^{\mathrm{L}}_k+R_k+\mathrm{OFF}\cdot2^b,
  \qquad R_0=R_K=\mathrm{OFF}=2^{J+6}.
\end{equation}
The offset makes the signed carry $R_k-\mathrm{OFF}$ a natural number. By \lean{carryChk_iff}, valid
carries exist if and only if $\sum_kd^{\mathrm{L}}_k2^{bk}=\sum_kd^{\mathrm{R}}_k2^{bk}$, and then they are
unique (\lean{carryChk_unique}) and lie in $(0,2^{J+7})$ (\lean{carryChk_range}). The mass bound $32$ keeps
every digit below $\mathrm{OFF}\cdot(2^b-1)$ (\lean{side_le_OFF}).
\item \emph{Comparisons: state chains.} A comparison $x_v\bowtie r$ has state wires $s_j\in\{0,1,2\}$
(less, equal, greater) passed from limb $\ell-1$ down to limb $0$; the cell at limb $i$ updates the state
lexicographically, and cell $0$ tests the relation (\lean{cmpChain_iff}, \lean{cmpChain_unique}).
\item \emph{Cells.} There is one cell per limb position $i$; it checks the range and pattern of limb $i$ of
every variable, that all copies agree, and the comparison transitions at $i$.
\end{itemize}
Every wire has exactly two ends. Index bookkeeping goes through a generic mixed-radix block layout
(module \lean{Blocks}) instead of hand-written case splits.

\subsection{Value and parameters}
For a description with $\mathrm{nV}$ variables, $\mathrm{nP}$ products, $\mathrm{nC}$ chains and
$\mathrm{nQ}$ comparisons let $c_A(D)=64(\mathrm{nV}+\mathrm{nP}+\mathrm{nQ}+1)(\mathrm{nC}+1)$ and
$d(n)=2^{\lfloor3\nu/4\rfloor+11}=O(n^{3/4})$.

\begin{theorem}[Parsimony; \lean{arithNet_card}, \lean{arithNet_value}]\label{thm:value}
Let $g\ge2$ and let $D$ be well formed at size $n$. The labellings accepted by every check are in bijection
with the solutions counted by $\mathrm{count}_n(D)$, so the value of the network is $\mathrm{count}_n(D)$.
\end{theorem}

\begin{proof}[Sketch]
\emph{Forward} (\lean{aRun_all_iff}): a solution $x$ labels every wire with its canonical value
$\mathrm{run}(x)$; the transform checks pass by construction, the pipelines compute the products
(\lean{pipeline_eq_mul} of~\cite{TulFac26}), so the carry checks pass if and only if every chain balances,
and the cells pass if and only if the patterns match and the comparisons hold. \emph{Backward}
(\lean{aEq_run}): an accepted labelling equals $\mathrm{run}$ of the values read off the limb wires,
stage by stage, with the carries and states forced by uniqueness.
\end{proof}

\begin{proposition}[\lean{arithSys_U_le}, \lean{arithSys_is}]\label{prop:params}
For $g\ge2$ the network has total $U\le2^{11}c_A(D)(n+2)$, and it is a network with
$P=2^{12}c_A(D)\bigl((n+2)/g+d(n)\bigr)$ and $B=1$.
\end{proposition}

\noindent The network is in \A's \lean{Net} format through our label systems~\cite{Tul26}, and its
tables are answered by polynomial-time procedures (\lean{tabAF_polyT}, \lean{arithNet_isNet}); no table is
written out. For the descriptions of \cref{sec:instances}, $c_A(D)\le1152$.

% =====================================================================
\section{Counted searches and queries}\label{sec:machines}
% =====================================================================

In this section we describe how counts become answers. The interface is a \emph{counting package}
\lean{CntPackage}~$c_0$: for every input word $w$ and query word $\pi$ a network whose value is
$Q(w,\pi)\in\N$, with $U\le c_0(|w|+2)$, $P\le c_0((|w|+2)/g+d(|w|))$, polynomial reader and procedures,
and a number of rounds $R(w)\le|w|$. \lean{ArithSpec.pkg_exists} builds a package with
$Q(w,\pi)=\mathrm{count}_{n(w)}(D(w,\pi))$ from any polynomial-time map $(w,\pi)\mapsto D(w,\pi)$ to
well-formed descriptions with $c_A\le c_D$, with $c_0=2^{12}c_D$.

\paragraph{Prefix search.} A \lean{SearchSpec} has a test for small inputs, an answer for them, a step
and a finish, all polynomial-time. The machine \lean{GR.searchM} starts from the empty word, and for
$R(w)$ rounds replaces $\pi$ by $\mathrm{step}(w,\pi,Q(w,\pi0))$; then it outputs $\mathrm{fin}(w,\pi)$.
With the step ``append $0$ if the count is positive, else $1$'', the word found is the least element of
the counted set (\lean{searchWith_eq_minOr}: if the oracle is positive on $\pi0$ exactly when some element of a
set $S\subseteq[0,2^n)$ has prefix $\pi0$, the result has value $\min S$, or $2^n-1$ if $S=\emptyset$).

\begin{theorem}[\lean{GR.searchM_TM}, \lean{GR.queryM_TM}]\label{thm:machines}
For every counting package and every search specification (resp.\ query specification with two query words
$q_1(w),q_2(w)$ and a finish $\mathrm{fin}(w,Q(w,q_1),Q(w,q_2))$) whose parts are polynomial-time and whose
outputs have length $O(|w|)$, there are a finite-alphabet Turing machine and $C$ that output the specified
answer on every $w$ within $\lfloor C\cdot2^{C|w|/\sqrt{\log_2(|w|+2)}}\rfloor$ steps.
\end{theorem}

\noindent Both are proved through one generic tail, \lean{GR.TM_of_runs}, on top of the round bound
\lean{FR.round_bound} and the level selection of the factoring note.

% =====================================================================
\section{Six descriptions}\label{sec:instances}
% =====================================================================

In this section we give each description and the pure counting identity that connects its count to
the answer. All identities are in pure Mathlib (modules \lean{ClassForms}, \lean{QRCount},
\lean{Lattice}, \lean{ArithInstances}).

\subsection{Reduced forms with two queries}
We follow Cox's convention~\cite[\S2]{Cox13}: a positive definite form $(a,b,c)$ of discriminant
$D=b^2-4ac<0$ is \emph{reduced} if $|b|\le a\le c$, and $b\ge0$ whenever $|b|=a$ or $a=c$.

\begin{definition}[\lean{hForm}]\label{def:hform}
For $D\in\Z$ and $N=|D|$, $\hF(D)$ is the number of $(b,a,c)\in[-N,N]\times[0,N]\times[0,N]$ with
$b^2-4ac=D$, $|b|\le a\le c$, and $(|b|=a\vee a=c)\Rightarrow b\ge0$.
\end{definition}

\noindent For $D<0$ the box imposes nothing (\lean{mem_redForms_iff}): a reduced form has $c\le N$. The
description of query~$j$ has variables $a,b,c$, products $b\cdot b$ and $a\cdot c$, the chain
$1\cdot(b\cdot b)+1\cdot N=4\cdot(a\cdot c)$ with $N$ a constant, and the comparisons
\[
  \text{query 1: } b\le a,\ a\le c;\qquad \text{query 2: } b>0,\ b<a,\ a<c .
\]
Here $b$ ranges over naturals. Query~1 counts the reduced forms with $b\ge0$; in them the tie-break is
automatic. A reduced form with $b<0$ has $|b|<a<c$, since equality would force $b\ge0$, so $b\mapsto-b$
maps these forms bijectively onto the solutions of query~2.

\begin{lemma}[\lean{hForm_split}, \lean{hDesc_count}]\label{lem:split}
For $0<N<B$, $\hF(-N)=h_{Q1}(B,N)+h_{Q2}(B,N)$, where $h_{Qj}(B,N)$ is the number of solutions of query~$j$
in $[0,B)^3$. The two descriptions count exactly these.
\end{lemma}

\noindent The machine (\lean{specE}) answers $1$ on $N=0$ (where $\hF(0)=1$ counts the degenerate triple
$(0,0,0)$), and otherwise asks the two queries and adds. In \cref{ex:23} the counts are $2$ and $1$.

\subsection{Square roots modulo any \texorpdfstring{$m$}{m}}
The input is \lean{encodePair}~$m$~$a$, two little-endian words padded to a common length $n$. The reader
takes $m$ and $a\bmod m$; inputs with $m<2$ are answered directly ($0$). The description has variables
$y,k,x$, products $y\cdot y$ and $k\cdot m$, the chain $y\cdot y=x+k\cdot m$, the comparisons $y<m$, $x<m$,
$x=a\bmod m$, and the pattern ``$y$ has the prefix $\pi0$''. Since $y<m$ and $x<m$ determine $k$ and $x$,
the count is the number of roots $y<m$ of $a$ with that prefix (\lean{S3A_eq}). The search keeps $0$ when the
count is positive; at the end, $y=\val(\pi)$ is output if it is a root and $m$ otherwise, which covers the
case of no roots, where the search returns $2^n-1\ge m$. For $m=7$, $a=2$ (roots $3,4$; $n=3$) the counts of
$0$, $00$, $010$ are $1,0,0$, and the result is $011=3$.

\subsection{The least non-residue from counts of squares}
For a prime $p$ we need, for each prefix, whether some non-residue $x<p$ has it. The same description
with the pattern on $x$ instead of $y$, and no condition $x=a$, counts the pairs $(y,x)$ with
$y^2\equiv x$ and $x$ in the prefix block. For an odd prime, $0$ has one square root and every nonzero
square has two, so this count is $S(\pi)=2\,Q(\pi)+z_0(\pi)$, where $Q(\pi)$ counts the nonzero squares in the
block and $z_0(\pi)=1$ if the block contains $0$ (\lean{S3B_eq_Qc}). The block has $r(\pi)$ elements below
$p$, so the number of non-residues in it is
\begin{equation}\label{eq:nr}
  \mathrm{NR}(\pi)=r(\pi)-z_0(\pi)-\bigl(S(\pi)-z_0(\pi)\bigr)/2\qquad\text{(\lean{NR_eq})},
\end{equation}
which the step computes from the count in polynomial time; it keeps $0$ when $\mathrm{NR}(\pi0)>0$.

\begin{example}[$p=7$]\label{ex:qnr}
The squares modulo $7$ are $0,1,2,4$, so the least non-residue is $3=011_2$. The search runs as follows.
\[
\begin{array}{l|ccc}
\text{query }\pi0 & 0 & 00 & 010\\\hline
\text{block} & \{0,1,2,3\} & \{0,1\} & \{2\}\\
S(\pi0) & 5 & 3 & 2\\
r,\ z_0 & 4,\ 1 & 2,\ 1 & 1,\ 0\\
\mathrm{NR}(\pi0) & 1 & 0 & 0\\
\text{new }\pi & 0 & 01 & 011
\end{array}
\]
\end{example}

\noindent For $p\le2$ the machine answers $p$ ($0$ and $1$ are squares modulo $2$). For $p$ composite the
formula may be wrong, but the finish outputs some $y\le p$ (\lean{specB_answer_le}). The pure
characterisation is \lean{lqnr_spec}: for a prime $p>2$, $\mathrm{lqnr}(p)<p$, it is not a square in
\lean{ZMod}~$p$, and every smaller residue is a square.

\subsection{Divisor sums and the circle}\label{sec:lattice}
These need one query each. The divisor sum is $\sum_{m\le x}\tau(m)=\#\{(a,b,s):ab+s=x,\ a,b>0\}$
(\lean{divQ_eq}); $\sum_{m\le x}\sigma(m)=\#\{(a,b,c,s):ab+s=x,\ b>0,\ 1\le c\le a\}$ (\lean{sigQ_eq}); and the
circle count is $4\,\#\{(a,b,s):a^2+b^2+s=x,\ a>0\}+1$, by rotating the four quadrants
(\lean{circ_eq}). Each is one product (two for the circle) and one chain with the constant $x$.

\begin{example}\label{ex:lattice}
For $x=10$: $\sum_{m\le10}\tau(m)=27$, $\sum_{m\le10}\sigma(m)=87$, and the circle count is $4\cdot9+1=37$.
\end{example}

% =====================================================================
\section{What the class-number theorem counts}\label{sec:honesty}
% =====================================================================

In this section we say exactly what Lean proves about class numbers and what is classical.

\paragraph{Forms, not ideals.} Lean proves that the machine outputs $\hF(-N)$, the number of Gauss reduced
forms of discriminant $-N$, and, for a fundamental discriminant, the number of \emph{primitive} reduced forms
(\lean{hForm_eq_card_primitive}, from \lean{primitive_of_fund}: no prime divides all of $a,b,c$ when
$b^2-4ac$ is fundamental). That number is the form class number $h(D)$, the order of Gauss's form class
group. That $h(D)$ equals the ideal class number of the imaginary quadratic field $\mathbb{Q}(\sqrt D)$
for fundamental $D$, and of the order of discriminant $D$ in general, is the classical correspondence of
Dirichlet and Dedekind~\cite[Thm.~7.7]{Cox13}, \cite{Bue89}. It is not in Mathlib at the commit we use:
Mathlib has \lean{NumberField.classNumber}, but no binary quadratic forms and no forms--ideals
correspondence. So the link to \lean{NumberField.classNumber} is proved on paper only. We define
fundamental discriminants locally (\lean{IsFundDisc}: $D\equiv1\pmod 4$ squarefree, or $D=4m$ with
$m\equiv2,3\pmod4$ squarefree), since Mathlib at that commit has no such definition.

\begin{remark}[Non-fundamental discriminants]\label{rem:nonfund}
If $N\equiv1,2\pmod4$ there are no forms, and $\hF(-N)=0$. Otherwise $\hF(-N)$ also counts imprimitive forms:
every form of discriminant $D$ is $f$ times a primitive form of discriminant $D/f^2$, and scaling preserves
reduction, so
\[
  \hF(D)=\sum_{f^2\mid D,\ D/f^2\equiv0,1\ (4)}h(D/f^2),\qquad
  h(D)=\sum_{f^2\mid D,\ D/f^2\equiv0,1\ (4)}\mu(f)\,\hF(D/f^2).
\]
For example $\hF(-12)=2$ counts $(1,0,3)$ and $(2,2,2)=2\cdot(1,1,1)$, while $h(-12)=1$. The inversion
recovers $h(D)$ from at most $\tau(N)$ runs of the machine once the square divisors of $N$ are known, and
these can be found by factoring $N$, also in time $\rate$~\cite{TulFac26}. This is on paper only (we checked
the decomposition numerically for $N<400$); the Lean theorem for non-fundamental $D$ is about $\hF$.
\end{remark}

\paragraph{Real quadratic fields are not covered.} For $D>0$ the reduced indefinite forms fall into
cycles under the reduction operator, and each class corresponds to a whole cycle, whose length is tied
to the regulator. Counting reduced forms therefore gives neither $h$ nor $hR$: one must count cycles, or
weight each form by its distance to its successor, which yields $hR$~\cite[\S4--5]{Sch82}. The
indicator ``this form is the first of its cycle'' is a condition over the whole cycle, not a local check
on $O(n)$ bits, and separating $h$ from the regulator $R$ is a baby-step giant-step computation, not a
count. We know of no description with few products for $h(D)$, $D>0$.

% =====================================================================
\section{The exponent}\label{sec:rate}
% =====================================================================

In this section we check that the descriptions give the rate. Every description of \cref{sec:instances}
has at most four variables, two products, one chain and three comparisons, so $c_A(D)\le1152=c_D$ for every
input. By \cref{prop:params}, $U\le2^{11}\cdot1152\,(n+2)$ and $gP\le2^{12}\cdot1152\,((n+2)+g\,d(n))$. At
the level $K$ with $K^{40}\le K_0^{40}(n+2)$ of~\cite[\S5]{TulFac26} we have $g\,d(n)=O(n)$, so $U,gP=O(n)$,
and \cref{imp:engine} gives one count in $2^{O(n/\sqrt{\log n})}$ steps. The searches ask at most $n$
counts and the class-number and lattice machines one or two; the overhead of the host lemma is
polynomial. In Lean this is \lean{FR.round_bound} with $c_0=2^{12}\cdot1152$, inside
\cref{thm:machines}. For (A) the network lives at $n=|w|/2$ while the bounds are stated in $|w|$; this
uses that $d$ is monotone (\lean{dg_mono}).

% =====================================================================
\section{What a reader must trust}\label{sec:statement}
% =====================================================================

A reader who wants to trust \cref{thm:main} must check the following definitions and nothing else of
ours. They are in the pure modules \lean{ClassForms}, \lean{QRCount} and \lean{Encode}; we render Unicode
symbols as mathematics (\verb|x.1| and \verb|x.2| are the projections of a pair).
\begin{alltt}\footnotesize
def hForm (D : \(\mathbb{Z}\)) : \(\mathbb{N}\) :=
  ((Finset.Icc (-(D.natAbs : \(\mathbb{Z}\))) D.natAbs \(\times\)\(\sp{s}\) Finset.Icc 0 (D.natAbs : \(\mathbb{Z}\)) \(\times\)\(\sp{s}\)
      Finset.Icc 0 (D.natAbs : \(\mathbb{Z}\))).filter
    fun t : \(\mathbb{Z}\) \(\times\) \(\mathbb{Z}\) \(\times\) \(\mathbb{Z}\) => t.1 * t.1 - 4 * t.2.1 * t.2.2 = D \(\wedge\) |t.1| \(\le\) t.2.1 \(\wedge\)
      t.2.1 \(\le\) t.2.2 \(\wedge\)
      ((|t.1| = t.2.1 \(\vee\) t.2.1 = t.2.2) \(\to\) 0 \(\le\) t.1)).card
def redForms (D : \(\mathbb{Z}\)) : Finset (\(\mathbb{Z}\) \(\times\) \(\mathbb{Z}\) \(\times\) \(\mathbb{Z}\)) := -- the same finset as in hForm
def IsFundDisc (D : \(\mathbb{Z}\)) : Prop :=
  (D % 4 = 1 \(\wedge\) Squarefree D) \(\vee\)
    (D % 4 = 0 \(\wedge\) (D / 4 % 4 = 2 \(\vee\) D / 4 % 4 = 3) \(\wedge\) Squarefree (D / 4))
def primRedForms (D : \(\mathbb{Z}\)) : Finset (\(\mathbb{Z}\) \(\times\) \(\mathbb{Z}\) \(\times\) \(\mathbb{Z}\)) :=
  (redForms D).filter fun t => Int.gcd t.2.1 (Int.gcd t.1 t.2.2) = 1
def minRoot (m a : \(\mathbb{N}\)) : \(\mathbb{N}\) :=
  Nat.find (\(\langle\)m, Or.inl rfl\(\rangle\) : \(\exists\) y, y = m \(\vee\) (y < m \(\wedge\) y * y % m = a % m))
def lqnr (p : \(\mathbb{N}\)) : \(\mathbb{N}\) :=
  Nat.find (\(\langle\)p, Or.inl rfl\(\rangle\) : \(\exists\) x, x = p \(\vee\) (x < p \(\wedge\) \(\forall\) y < p, y * y % p \(\ne\) x))
def padR (n : \(\mathbb{N}\)) (l : List Bool) : List Bool := l ++ List.replicate (n - l.length) false
def encodePair (m a : \(\mathbb{N}\)) : List Bool :=
  padR (max (encodeNat m).length (encodeNat a).length) (encodeNat m) ++
    padR (max (encodeNat m).length (encodeNat a).length) (encodeNat a)
\end{alltt}
A triple is \texttt{(b, a, c)}, so \texttt{t.1} is the middle coefficient. The final statements, verbatim
from \lean{Final} up to line breaks, except that we abbreviate the common bound as $\mathrm{T}(x)$, for the input word $x$,
and write \texttt{Outputs} for \lean{GraphProblems.Outputs}. In the source, $\mathrm{T}(x)$ is
\begin{alltt}\footnotesize
\(\lfloor\)(C : \(\mathbb{R}\)) * 2 ^ ((C : \(\mathbb{R}\)) * x.length /
    Real.sqrt (Real.logb 2 (x.length + 2)))\(\rfloor\sb{+}\)
\end{alltt}
that is, $\lfloor C\cdot2^{C|x|/\sqrt{\log_2(|x|+2)}}\rfloor$. The statements are:
\begin{alltt}\footnotesize
theorem hForm_TM_rate :
  \(\exists\) (M : Turing.TM2ComputableAux Bool Bool) (C : \(\mathbb{N}\)), (\(\forall\) i, Finite (M.tm.\(\Gamma\) i)) \(\wedge\)
  \(\forall\) N : \(\mathbb{N}\), Outputs M (encodeNat N) (encodeNat (hForm (-(N : \(\mathbb{Z}\))))) \(\mathrm{T}\)(encodeNat N)
theorem hFormPrim_TM_rate : \(\exists\) M C, (\(\forall\) i, Finite (M.tm.\(\Gamma\) i)) \(\wedge\)
  \(\forall\) N : \(\mathbb{N}\), IsFundDisc (-(N : \(\mathbb{Z}\))) \(\to\) 0 < N \(\to\)
    Outputs M (encodeNat N) (encodeNat (primRedForms (-(N : \(\mathbb{Z}\)))).card) \(\mathrm{T}\)(encodeNat N)
theorem sqrtMod_TM_rate : \(\exists\) M C, (\(\forall\) i, Finite (M.tm.\(\Gamma\) i)) \(\wedge\)
  \(\forall\) m a : \(\mathbb{N}\), Outputs M (encodePair m a) (encodeNat (minRoot m a)) \(\mathrm{T}\)(encodePair m a)
theorem lqnr_TM_rate : \(\exists\) M C, (\(\forall\) i, Finite (M.tm.\(\Gamma\) i)) \(\wedge\)
  (\(\forall\) p : \(\mathbb{N}\), p.Prime \(\to\)
     Outputs M (encodeNat p) (encodeNat (lqnr p)) \(\mathrm{T}\)(encodeNat p)) \(\wedge\)
  (\(\forall\) p : \(\mathbb{N}\), \(\exists\) y \(\le\) p, Outputs M (encodeNat p) (encodeNat y) \(\mathrm{T}\)(encodeNat p))
theorem divisorSum_TM_rate : \(\exists\) M C, (\(\forall\) i, Finite (M.tm.\(\Gamma\) i)) \(\wedge\) \(\forall\) x : \(\mathbb{N}\),
  Outputs M (encodeNat x) (encodeNat (\(\sum\) m \(\in\) Finset.Icc 1 x, (Nat.divisors m).card))
    \(\mathrm{T}\)(encodeNat x)
theorem sigmaSum_TM_rate : \(\exists\) M C, (\(\forall\) i, Finite (M.tm.\(\Gamma\) i)) \(\wedge\) \(\forall\) x : \(\mathbb{N}\),
  Outputs M (encodeNat x) (encodeNat (\(\sum\) m \(\in\) Finset.Icc 1 x, \(\sum\) d \(\in\) Nat.divisors m, d))
    \(\mathrm{T}\)(encodeNat x)
theorem circle_TM_rate : \(\exists\) M C, (\(\forall\) i, Finite (M.tm.\(\Gamma\) i)) \(\wedge\) \(\forall\) x : \(\mathbb{N}\),
  Outputs M (encodeNat x)
    (encodeNat (((Finset.Icc (-(x : \(\mathbb{Z}\))) x \(\times\)\(\sp{s}\) Finset.Icc (-(x : \(\mathbb{Z}\))) x).filter
      fun q : \(\mathbb{Z}\) \(\times\) \(\mathbb{Z}\) => q.1 ^ 2 + q.2 ^ 2 \(\le\) x).card)) \(\mathrm{T}\)(encodeNat x)
\end{alltt}
In the source, each statement spells out the bound and writes \texttt{Computability.encodeNat}; the
types of \texttt{M} and \texttt{C} are as in the first. \texttt{Outputs} is the one-line definition of \A's
statement file, which we also trust:
\begin{alltt}\footnotesize
def Outputs (M : TM2ComputableAux Bool Bool) (x y : List Bool) (T : \(\mathbb{N}\)) : Prop :=
  Nonempty (TM2OutputsInTime M.tm (x.map M.inputAlphabet.invFun)
    (some (y.map M.outputAlphabet.invFun)) T)
\end{alltt}
Behind these are Mathlib's \lean{TM2OutputsInTime}, \lean{Computability.encodeNat} (little-endian
binary, the empty word for $0$), \lean{Nat.find}, \lean{Nat.divisors}, \lean{Finset.Icc},
\lean{Squarefree}, \lean{Int.gcd}, \lean{Real.logb} and \lean{Real.sqrt}, from Mathlib \texttt{0df444a}.
Integer division and remainder are Lean's (\lean{Int.emod} is non-negative), which in \lean{IsFundDisc}
gives the usual definition for $D<0$. For a quick check of the definitions, Lean evaluates
$\hF(-3)=\hF(-4)=\hF(-7)=\hF(-8)=1$, $\hF(-15)=\hF(-20)=2$ and $\mathrm{lqnr}(7,11,13,17,41)=3,2,2,3,3$ by
\lean{decide}.

% =====================================================================
\section{The Lean development}\label{sec:formal}
% =====================================================================

In this section we list the modules and state the trusted base. The development is in the public
repository~\cite{TulsyanRepo26}: directories \texttt{lean/EthFollowup/ModP} (pure Mathlib; the files
compile with Lean v4.33.1 and v4.34.1) and \texttt{lean-bridge-anthropic/bridge/EthFollowup/ModP} (on top
of \A's development). \Cref{tab:modules} lists the modules.

\begin{table}[ht]
\centering\small
\begin{tabular}{@{}p{0.36\textwidth}p{0.50\textwidth}r@{}}
\toprule
module & content & lines\\
\midrule
\lean{SignedCarry} & signed carry chains: existence iff sum, uniqueness, range & 417\\
\lean{Compare} & comparison state chains & 213\\
\lean{MinSearch} & prefix search finds the minimum & 164\\
\lean{QRCount} & \lean{minRoot}, \lean{lqnr}, root and square counts, \eqref{eq:nr} & 509\\
\lean{Encode} & \lean{encodePair} and its halves & 146\\
\lean{ClassForms} & \lean{hForm}, the two-query split, primitivity & 277\\
\lean{Lattice} & divisor sums and the circle as solution counts & 293\\
\lean{ArithDesc} & descriptions and their count & 192\\
\lean{ArithInstances} & the six descriptions and their counts & 427\\
\midrule
\lean{Blocks} & mixed-radix block layouts & 248\\
\lean{ArithBounds}, \lean{ArithParams} & digit bounds; wires, vertices, checks, canonical run & 957\\
\lean{ArithNetDefs}, \lean{ArithNetParsimony}, \lean{ArithNetCount}, \lean{ArithNet} & the label system, parsimony, $U$ and $P$ & 947\\
\lean{ArithEnc}, \lean{ArithReaderPoly}, \lean{ArithProcs}, \lean{ArithPoly} & encodings, procedures, polynomial time & 1647\\
\lean{CntPackage}, \lean{GRate} & counting packages; search and query machines & 878\\
\lean{ArithPackage}, \lean{ArithPkgs} & packages from descriptions; the six packages & 589\\
\lean{SpecSq}, \lean{SpecE}, \lean{SpecLattice} & the search and query specifications & 616\\
\lean{Final} & the seven final theorems & 112\\
\midrule
total & & 8632\\
\bottomrule
\end{tabular}
\caption{The modules, with line counts (\texttt{wc -l}). The first nine are pure Mathlib. The sanity module
\texttt{GRateMinFac} (77 lines, re-deriving the factoring theorem through \texttt{GR.searchM\_TM}) is not
imported by \texttt{Final}. The development reuses, unchanged, the factoring modules of~\cite{TulFac26}
(\texttt{Params}, \texttt{FermatBridge}, \texttt{MulProcs}, \texttt{MulPackage}, \texttt{Level40}, \texttt{Rate},
\texttt{Search}) and the bridge modules of~\cite{Tul26}.}\label{tab:modules}
\end{table}

\paragraph{Trusted base and axioms.} The theorems depend only on \lean{propext}, \lean{Classical.choice}
and \lean{Quot.sound}.
A fresh \lean{leanchecker} replay of the dependency closure of \lean{EthFollowup.ModP.Final}, the module
holding the seven final theorems, has passed.
What must be trusted is the Lean kernel, the Mathlib definitions named in \cref{sec:statement}, our
definitions quoted there, and \A's definition of \texttt{Outputs}. \A's development, Mathlib's proofs and
our other modules are checked by the kernel; they are part of the checked closure, not of the trusted base.

\paragraph{One design point.} Everything polynomial-time must be polynomial on \emph{every} context, not
only on those the reader produces. Our first context stored the number of variables and vertices in binary
and looped over them, which is exponential on adversarial contexts; the context now stores them in unary
(\lean{ArithProcs}).

% =====================================================================
\section{Related work}\label{sec:related}
% =====================================================================

We organise related work by idea: non-residues and square roots, class numbers, lattice-point counts,
the counting engines applied to number theory, and formal proofs.

\paragraph{Non-residues and square roots.} Burgess's bound $p^{1/(4\sqrt e)+\varepsilon}$~\cite{Bur57},
from his character-sum method~\cite{Bur62}, is the best unconditional bound for the least non-residue.
Under GRH it is $O(\log^2p)$ by Ankeny~\cite{Ank52}, with explicit constants by Bach~\cite{Bac90} and
Lamzouri, Li and Soundararajan~\cite{LLS15}. Given a non-residue, square roots modulo $p$ take polynomial
time by Tonelli~\cite{Ton91} and Shanks~\cite{Sha72}, or Cipolla~\cite{Cip03}; Schoof~\cite{Sch85}
computes $\sqrt a\bmod p$ in polynomial time for fixed $a$. Finding a non-residue and taking square roots
modulo $p$ are deterministic polynomial-time equivalent (folklore;~\cite{BS96}). For composite $m$, taking
square roots is randomised-equivalent to factoring $m$~\cite{Rab79}, so the classical deterministic route
first factors $m$, in $m^{1/5+o(1)}$~\cite{Har21}; our factoring note~\cite{TulFac26} would give a route in time $\rate$ as well, by factoring and then using~(iv), which we have not formalised. In contrast, we compute the \emph{least} root and the
\emph{least} non-residue directly by counting, for every $m$, conditionally on \A's development.

\paragraph{Class numbers.} Counting reduced forms is Gauss's method and is used to tabulate class
numbers~\cite[\S2]{Sch82}, \cite{Bue89}; the best proven deterministic unconditional bound we know is
$|D|^{1/2+o(1)}$~\cite[\S5]{Len92}. Under GRH, Shanks's method gives $|D|^{1/5+o(1)}$
deterministically~\cite{Len82,Sch82}, \cite[\S5.4]{Coh93}, and Hafner and McCurley's Las Vegas algorithm
runs in expected time $L(|D|)^{\sqrt2+o(1)}$, $L(x)=\exp\sqrt{\log x\log\log x}$~\cite{HM89}; its
analysis assumes GRH~\cite[\S5]{Len92}. Mosunov and Jacobson~\cite{MJ16} tabulate all class groups with
$|D|<2^{40}$ unconditionally, which is a different, amortised problem. In contrast, our algorithm is
deterministic and unconditional except for \A's development, and it is $|D|^{o(1)}$; but it computes only
the number, not the group.

\paragraph{Lattice-point counts.} Sladkey~\cite{Sla12} computes $\sum_{m\le x}\tau(m)$ in $O(x^{1/3})$ steps,
and Tao, Croot and Helfgott~\cite[\S2.1]{TCH12} sketch an $x^{1/3+o(1)}$ algorithm; Li~\cite{Li26} reports
$x^{9/28+o(1)}$ expected time. For the circle count we found no dedicated reference; we expect the
same convex-hull techniques to give $x^{1/3+o(1)}$. In contrast, we obtain $x^{o(1)}$, at the price of galactic
constants.

\paragraph{Counting engines applied to number theory.} Rudich and Rousseau~\cite{RR26} introduced prefix
search by divisor counting and proved factoring in $N^{o(1)}$ in Lean; Shunia~\cite{Shu26} proved the rate
$\rate$ on paper, and~\cite{IF26} and our factoring note~\cite{TulFac26} proved it in Lean. Closest to our
work is the factoring note, whose multiplication network and level selection we reuse. In contrast, we
replace its single network by a library for any few-product system and apply it to problems other than
factoring. As of 4~November 2026 we found no derivation of the bounds of \cref{thm:main} from either ETH
refutation. We searched the web for each problem together with ``ETH'', ``$\#$SAT'' and ``2026''; searched
GitHub repositories for the problem names; checked the Hexagon records and repositories of the October
and November 2026 follow-ups known to us (\cite{RR26,Shu26,IF26} and the fork
\repo{orda-aska/integer-factoring}), which treat factoring only; and read our survey of
targets~\cite{TulsyanRepo26}. Hexagon's listing is behind a login, so this negative is weak.

\paragraph{Formal proofs.} We use Lean~4~\cite{Lean4} and Mathlib~\cite{Mathlib20}. Mathlib defines the class
number of a number field (\lean{NumberField.classNumber}) and proves its finiteness, but has no theory of
binary quadratic forms. In contrast, our statements are about forms, and the bridge to ideals is the
classical theorem cited in \cref{sec:honesty}.

% =====================================================================
\section{Limitations}\label{sec:limitations}
% =====================================================================

\begin{itemize}
\item \emph{The results stand or fall with \A's development.} Every theorem rests on \A's Lean development
of its claimed refutation of the Exponential Time Hypothesis, and the labs' claims are under review. Lean
checks that development as part of our proof, and our statements use only Mathlib, our quoted definitions
and \A's \texttt{Outputs}. We do need the development, built from source at the pinned commit, to be a
sound Lean development, and the toolchain to be sound.
\item \emph{The two number-theoretic headlines share one source of strength.} Finding a non-residue and
taking square roots modulo $p$ are classically polynomial-time equivalent~\cite{BS96}, so (iii) for prime $m$
and (iv) do not give two independent improvements. Both, and (i) and (v), come from the same engine.
\item \emph{The class number is the form class number.} Its equality with the ideal class number is
classical but not formalised; for non-fundamental $D$ the theorem is about $\hF$
(\cref{rem:nonfund}); the class group and real quadratic fields are not covered.
\item \emph{The algorithms are galactic.} \A's levels start at $K_0=2^{x_0}$ with $x_0\ge18$, the level
needs $K^{40}\lesssim n$, and the constants $c_0=2^{12}\cdot1152$ and $C$ are astronomical. The machines are
of no practical use, and they beat the classical bounds only for astronomically large inputs.
\item \emph{The machines use exponential space.}
\item \emph{No human has checked the work.} The Lean code and this text were written with AI assistance.
No human expert has yet checked that \cref{sec:statement} lists everything a reader must trust.
\end{itemize}

\section*{Acknowledgements}
This work was done with AI assistance (Claude Code, Anthropic): the constructions, the Lean development,
the literature and citation checks and the text were produced in sessions with Claude Code. The Lean
theorems are checked by the Lean kernel; everything else awaits human expert review. Responsibility for any
errors rests with the author. This note continues the project of~\cite{Tul26,TulFac26}.

\begin{thebibliography}{99}

\bibitem{Ank52}
N.~C.~Ankeny, \emph{The least quadratic non residue}, Ann. of Math. (2) \textbf{55} (1952), 65--72.

\bibitem{Anthropic26}
Claude (Anthropic), \emph{A $2^{o(n)}$ time algorithm for {3-SAT}}, preprint dated 26~October 2026.
Paper and Lean formalisation in the repository \url{https://github.com/anthropics/formal-math}, folder
\texttt{eth}; we use commit \texttt{16a5ce41304fcdaea4e0be7c132e1c361dc26da5}.

\bibitem{Bac90}
E.~Bach, \emph{Explicit bounds for primality testing and related problems}, Math. Comp. \textbf{55}
(1990), no.~191, 355--380.

\bibitem{BS96}
E.~Bach and J.~Shallit, \emph{Algorithmic Number Theory, Vol.~1: Efficient Algorithms}, MIT Press,
Cambridge, MA, 1996.

\bibitem{Bue89}
D.~A.~Buell, \emph{Binary Quadratic Forms: Classical Theory and Modern Computations}, Springer-Verlag,
New York, 1989.

\bibitem{Bur57}
D.~A.~Burgess, \emph{The distribution of quadratic residues and non-residues}, Mathematika \textbf{4}
(1957), 106--112.

\bibitem{Bur62}
D.~A.~Burgess, \emph{On character sums and primitive roots}, Proc. London Math. Soc. (3) \textbf{12}
(1962), 179--192.

\bibitem{Cip03}
M.~Cipolla, \emph{Un metodo per la risoluzione della congruenza di secondo grado}, Rend. Accad. Sci.
Fis. Mat. Napoli (3) \textbf{9} (1903), 154--163.

\bibitem{Coh93}
H.~Cohen, \emph{A Course in Computational Algebraic Number Theory}, Graduate Texts in Mathematics, vol.~138,
Springer-Verlag, Berlin, 1993.

\bibitem{Cox13}
D.~A.~Cox, \emph{Primes of the Form $x^2+ny^2$: Fermat, Class Field Theory, and Complex Multiplication},
2nd ed., John Wiley \& Sons, Hoboken, NJ, 2013.

\bibitem{Har21}
D.~Harvey, \emph{An exponent one-fifth algorithm for deterministic integer factorisation}, Math. Comp.
\textbf{90} (2021), no.~332, 2937--2950.

\bibitem{HM89}
J.~L.~Hafner and K.~S.~McCurley, \emph{A rigorous subexponential algorithm for computation of class
groups}, J. Amer. Math. Soc. \textbf{2} (1989), no.~4, 837--850.

\bibitem{IF26}
1337isnot1337 (GitHub user) and contributors, \emph{integer-factoring: deterministic integer factoring
bounds}, public repository \url{https://github.com/1337isnot1337/integer-factoring}; Lean formalisation in
\texttt{formal-appendix/} (pull request \#20, merged 30~October 2026).

\bibitem{LLS15}
Y.~Lamzouri, X.~Li, and K.~Soundararajan, \emph{Conditional bounds for the least quadratic non-residue and
related problems}, Math. Comp. \textbf{84} (2015), no.~295, 2391--2412; corrigendum, Math. Comp. \textbf{86}
(2017), 2551--2554.

\bibitem{Lean4}
L.~de Moura and S.~Ullrich, \emph{The {L}ean 4 theorem prover and programming language}, Automated
Deduction -- CADE 28, Lecture Notes in Comput. Sci., vol. 12699, Springer, 2021, pp.~625--635.

\bibitem{Len82}
H.~W.~Lenstra, Jr., \emph{On the calculation of regulators and class numbers of quadratic fields},
Journ\'ees Arithm\'etiques 1980 (J.~V.~Armitage, ed.), London Math. Soc. Lecture Note Ser., vol.~56,
Cambridge Univ. Press, 1982, pp.~123--150.

\bibitem{Len92}
H.~W.~Lenstra, Jr., \emph{Algorithms in algebraic number theory}, Bull. Amer. Math. Soc. (N.S.)
\textbf{26} (1992), no.~2, 211--244.

\bibitem{Li26}
D.~Li, \emph{Computing the divisor summatory function in $N^{9/28+o(1)}$}, arXiv:2609.29949, September 2026.

\bibitem{Mathlib20}
The mathlib Community, \emph{The {L}ean mathematical library}, Proc. 9th ACM SIGPLAN International
Conference on Certified Programs and Proofs (CPP 2020), 2020, pp.~367--381.

\bibitem{MJ16}
A.~S.~Mosunov and M.~J.~Jacobson, Jr., \emph{Unconditional class group tabulation of imaginary quadratic
fields to $|\Delta|<2^{40}$}, Math. Comp. \textbf{85} (2016), no.~300, 1983--2009.

\bibitem{OpenAI26}
OpenAI, \emph{The {E}xponential {T}ime {H}ypothesis is false}, preprint dated 17~October 2026, revised
25~October 2026. Paper and Lean formalisation in the repository \url{https://github.com/openai/math}.

\bibitem{Rab79}
M.~O.~Rabin, \emph{Digitalized signatures and public-key functions as intractable as factorization},
Technical Report MIT/LCS/TR-212, MIT Laboratory for Computer Science, January 1979.

\bibitem{RR26}
I.~Rudich and L.-M.~Rousseau, \emph{Deterministic factoring in time $N^{o(1)}$}, Hexagon preprint
hexagon:2610.00600v1, published 29~October 2026, \url{https://hexagonmath.org/2610.00600v1}.

\bibitem{Sch82}
R.~J.~Schoof, \emph{Quadratic fields and factorization}, Computational Methods in Number Theory, Part~II
(H.~W.~Lenstra, Jr. and R.~Tijdeman, eds.), Mathematical Centre Tracts, vol.~155, Mathematisch Centrum,
Amsterdam, 1982, pp.~235--286.

\bibitem{Sch85}
R.~Schoof, \emph{Elliptic curves over finite fields and the computation of square roots mod $p$}, Math.
Comp. \textbf{44} (1985), no.~170, 483--494.

\bibitem{Sha72}
D.~Shanks, \emph{Five number-theoretic algorithms}, Proc. Second Manitoba Conference on Numerical
Mathematics (Winnipeg, 1972), Congressus Numerantium VII, Utilitas Math., 1973, pp.~51--70.

\bibitem{Shu26}
J.~Shunia, \emph{An improved subexponential bound for deterministic integer factorization}, Hexagon
preprint hexagon:2610.00711v1, published 2~November 2026, \url{https://hexagonmath.org/2610.00711v1}.

\bibitem{Sla12}
R.~Sladkey, \emph{A successive approximation algorithm for computing the divisor summatory function},
arXiv:1206.3369, June 2012.

\bibitem{TCH12}
T.~Tao, E.~Croot III, and H.~Helfgott, \emph{Deterministic methods to find primes}, Math. Comp. \textbf{81}
(2012), no.~278, 1233--1246.

\bibitem{Ton91}
A.~Tonelli, \emph{Bemerkung \"uber die Aufl\"osung quadratischer Congruenzen}, Nachr. K\"onigl. Ges. Wiss.
G\"ottingen (1891), 344--346.

\bibitem{Tul26}
V.~Tulsyan, \emph{Counting {H}amiltonian cycles in sparse graphs in subexponential time, assuming the
new $\#${SAT} algorithms}, Hexagon preprint hexagon:2610.00672v1, October 2026; source and Lean
development at \url{https://github.com/xennygrimmato/eth-hamiltonicity}.

\bibitem{TulFac26}
V.~Tulsyan, \emph{An independent machine-checked proof of deterministic factoring in time
$2^{O(n/\sqrt{\log n})}$, via the new $\#${SAT} algorithms}, preprint, November 2026; source in the
directory \texttt{factoring-note/} of~\cite{TulsyanRepo26}.

\bibitem{TulsyanRepo26}
V.~Tulsyan, \emph{eth-hamiltonicity: {L}ean 4 formalisation and research notes}, public repository
\url{https://github.com/xennygrimmato/eth-hamiltonicity}, 2026; Lean sources as of commit \texttt{b9ea3f1}.

\end{thebibliography}

% =====================================================================
% CITATION CHECK (4 November 2026). Each entry was checked by web search or by fetching the source:
%  Ank52  Ann. of Math. (2) 55 (1952) 65-72, DOI 10.2307/1969420; GRH bound O(log^2 p) per citing works
%         (arXiv 2411.05095, MaRDI record).
%  Bur57  Mathematika 4 (1957) 106-112: found only in reference lists (no abstract online); exponent
%         1/(4 sqrt e) as in the standard literature (e.g. arXiv 2106.00544 cites it for this bound;
%         that preprint itself is NOT cited, unrefereed).
%  Bur62  Proc. LMS (3) 12 (1962) 179-192; mathnet.ru record of the Russian translation.
%  Bac90  Math. Comp. 55 (1990) no. 191, 355-380; AMS page, DOI 10.1090/S0025-5718-1990-1023756-8.
%  LLS15  Math. Comp. 84 (2015) no. 295, 2391-2412, arXiv 1309.3595; corrigendum Math. Comp. 86 (2017)
%         2551-2554 (AMS listing).
%  Ton91, Cip03, Sha72  Bibliographic data as in the standard secondary sources (Wikipedia
%         "Tonelli-Shanks algorithm", Dickson's History); Shanks's proceedings appeared 1973
%         (conference 1972). Originals not fetched.
%  Sch85  Math. Comp. 44 (1985) no. 170, 483-494; AMS page; square roots of fixed a in poly time
%         confirmed from the PDF (pages.cs.wisc.edu copy).
%  BS96   MIT Press 1996; equivalence of QNR finding and square roots attributed to it in lecture notes
%         (J. Booher, sqr_qnr.pdf) and in Sze, arXiv 0812.2591. Book itself not fetched.
%  Rab79  MIT/LCS/TR-212, January 1979; MIT DSpace handle 1721.1/149499.
%  Len92  Bull. AMS 26 (1992) no. 2, 211-244, arXiv math/9204234; full text read: Theorem 5.5 and
%         Problem 5.6 ("For quadratic fields the answer is affirmative": deterministic |D|^{1/2+o(1)}),
%         Problem 5.8 (GRH, L^{O(1)}) "For imaginary quadratic fields Problem 5.8 has been solved [28]"
%         where [28] = Hafner-McCurley.
%  HM89   J. AMS 2 (1989) no. 4, 837-850, DOI 10.1090/S0894-0347-1989-1002631-0. Abstract (geodesic.mathdoc)
%         says "rigorous proof ... Las Vegas ... L(d)^{sqrt2+o(1)}" and does not name a hypothesis;
%         Lenstra 1992 sec. 5 and de Boer-Pellet-Mary-Wesolowski (arXiv 2512.01588: "Assuming ERH, they
%         prove that the expected running time is L(1/2, sqrt2)") state that the analysis assumes
%         GRH/ERH. The AMS full text was not reachable (bot check). We state GRH, citing Lenstra.
%  Len82  Journees Arithmetiques 1980 (Exeter), LMS Lecture Note Ser. 56, CUP 1982, 123-150 (Springer
%         reference lists, Lenstra's publication list).
%  Sch82  Math. Centre Tracts 155 (1982) 235-286; full PDF read (mat.uniroma2.it/~schoof/mcpaper.pdf):
%         sec. 2 counting reduced forms for tables, sec. 3 Shanks's algorithm with a GRH derivation of the
%         order (X ~ |D|^alpha, alpha = 1/5), sec. 4-5 real quadratic fields: "one does not simply count
%         the forms, but one sums their distances to their successors".
%  Coh93  GTM 138, Springer 1993; table of contents: 5.3 class numbers of imaginary quadratic fields
%         (reduced forms), 5.4 class groups via Shanks's baby-step giant-step. The GRH 1/5 statement in
%         5.4 was not re-read (book not fetched); cited for the method.
%  "Buchmann |D|^{1/4} under GRH" (from the brief): NOT confirmed. Searches found |D|^{1/4+eps}
%         unconditional-probabilistic / |D|^{1/5+eps} under ERH attributions without a clear source,
%         so we do not cite a |D|^{1/4} bound.
%  MJ16   Mosunov, Jacobson, Math. Comp. 85 (2016) no. 300, 1983-2009, arXiv 1502.07953 (author order
%         Mosunov first, per AMS).
%  Sla12  arXiv 1206.3369 (15 June 2012), abstract: O(n^{1/3}) time, O(log n) space.
%  TCH12  Math. Comp. 81 (2012) no. 278, 1233-1246, arXiv 1009.3956; PDF read: sec. 2.1 "one can in
%         fact compute sum tau(n) in O(x^{1/3+o(1)}) time" (sketch).
%  Li26   arXiv 2609.29949 (D. Li, 24 Sep 2026): N^{9/28+o(1)} "expected time" per abstract.
%  Cox13  Wiley 2nd ed. 2013; Thm 7.7 (forms <-> ideal classes of an order) confirmed via citing papers
%         (arXiv 2308.11250 cites Thm 7.7(ii),(iii)). Bue89 Springer 1989 (standard; not fetched).
%  RR26, Shu26, IF26, Anthropic26, OpenAI26, Tul26, Lean4, Mathlib20: as audited in the factoring note
%         (factoring-note/main.tex, citation check of 4 Nov 2026). TulFac26: our own note, unpublished.
%  TulsyanRepo26 commit: git log -1 --format=%h -- lean-bridge-anthropic/bridge/EthFollowup/ModP/Final.lean
%         = b9ea3f1.
%  Novelty searches (4 Nov 2026): web searches "class number ... 2026 ... #SAT / ETH", "least quadratic
%         non-residue deterministic 2^{o(n)} 2026", "hexagonmath.org ... square roots / class number /
%         divisor summatory"; gh search repos for the problem names: nothing relevant found.
%  Mathlib 0df444a: NumberField.classNumber exists (Mathlib/NumberTheory/NumberField/ClassNumber.lean);
%         no binary quadratic forms, no FundamentalDiscriminant.
% =====================================================================

\end{document}
